\begin{abstract}
自然堆积场景下的骨料视觉筛分并不等价于“识别图像中的石子”。标准机械筛分以筛孔通过行为和逐级称重定义粒径级配，而单目图像直接获得的是二维可见轮廓及颗粒数量；因此，实例边界、像素尺度和统计口径之间存在连续的语义鸿沟。本文面向混凝土骨料颗粒智能筛分赛题，构建“实例感知--物理测量--筛分语义映射--任务输出”的分层方案：采用 ImageGrains 2.0 / Cellpose-SAM 预训练模型恢复颗粒实例，通过现场尺度标定获得毫米级长短轴、面积和形貌特征；进一步以短轴与等效直径构造可校准的筛分等效粒径，并用 $w=d^{\gamma}$ 将颗粒数量分布转换为质量相关分布，从而输出标准筛级质量占比及 $D_{10}/D_{50}/D_{90}$。同一实例几何还用于投影形貌分类和超限异常筛查。

当前三张自然堆积演示图像已形成从分割、几何测量到级配报告的完整闭环，共检测 3858 个实例。以 $\gamma=3$ 的未校准质量代理为例，三张图的数量型 $D_{50}$ 分别为 6.0、7.6 和 5.2~mm，而质量代理 $D_{50}$ 移动至 23.5、27.5 和 14.1~mm，说明从“按个数统计”到“按质量相关权重统计”会实质改变级配结论。仓库中的 27 项竞赛层合成测试已验证筛分映射、加权分位数、形貌、异常与报告逻辑的一致性；历史算法时序约为 27~s/张。需要强调的是，当前演示样本尚缺同批机械筛分逐筛称量真值，因此本文不将上述数值表述为筛分准确率。正式物理精度将通过“同批图像预测--机械筛分真值”配对实验校准并冻结 $\boldsymbol{\theta}$ 与 $\gamma$，再在独立批次上报告逐筛误差和 $D_p$ 误差。该方案的核心价值在于把预训练视觉能力与显式物理/统计校准结合，使最终输出既可自动化计算，也能追溯到颗粒实例与标准筛分语义。

\textbf{关键词：} 混凝土骨料；实例分割；Cellpose-SAM；粒径级配；机械筛分；质量加权
\end{abstract}
